We have an $r$-uniform hypergraph $G=(V,E)$, meaning each edge $e\in E$ satisfies $|e|=r$.

The hypothesis ``any two edges have a non-empty intersection'' means $G$ is \emph{intersecting}:
\[
\forall e,f\in E,\ e\neq f \implies e\cap f\neq \emptyset.
\]

\textbf{Ambiguity.} The text says ``chromatic number $3$'' but then parenthetically defines only the existence of a proper $3$-colouring (i.e. $\chi(G)\le 3$). Standard convention for ``chromatic number $3$'' is $\chi(G)=3$ (3-colourable but not 2-colourable).

\textbf{Minimal corrected interpretation.} I will interpret the intended condition as
\[
\chi(G)=3.
\]
I will also track what is true under the literal weaker assumption $\chi(G)\le 3$.

The problem asks:

(Q1) Is it true that every intersecting $r$-uniform hypergraph $G$ with $\chi(G)=3$ must satisfy $|V|=O(r^2)$?

(Q2) Is it true that there exists an absolute constant $c>0$ such that every such $G$ has two edges $e,f$ with $|e\cap f|\ge c r$?

(Q3) Besides the Fano plane example at $r=3$, are there intersecting $r$-uniform $\chi=3$ hypergraphs with \emph{no} pair of edges meeting in $r-1$ vertices?
